\documentclass[12pt,a4paper]{article}

\usepackage{styles/couleurs}
\usepackage{styles/mise_en_page}
\usepackage{styles/boites}
\usepackage{styles/maths}
\usepackage{styles/figures}

\usepackage{enumitem}
\usepackage{array}
\usepackage{tabularx}
\usepackage{amssymb}
\usepackage{tikz}

\setlist[enumerate]{leftmargin=*,itemsep=0.35em,topsep=0.35em}

\newcommand{\lignes}[1]{%
  \begin{tcolorbox}[colback=white,colframe=NeutreC!40,boxrule=0.5pt,arc=1.5mm,left=3mm,right=3mm,top=2mm,bottom=1mm]
    \foreach \i in {1,...,#1}{\noindent\dotfill\par\vspace{0.45em}}
  \end{tcolorbox}
}

\newcommand{\blocdonnees}[1]{%
  \begin{tcolorbox}[mathsbox,colback=NeutreBg,colframe=NeutreC!45,boxrule=0.6pt,arc=2mm]
  #1
  \end{tcolorbox}
}

\begin{document}

\begin{center}
  {\sffamily\LARGE\bfseries\color{TitreC}Séquence 1 -- Statistiques}\\[0.2em]
  {\sffamily\large\bfseries\color{SavoirC}Activité 4 -- Les grands lecteurs}
\end{center}

Deux classes du collège ont répondu à la question :

\begin{center}
\textit{« Combien de livres avez-vous empruntés durant les douze derniers mois ? »}
\end{center}

\blocdonnees{
\textbf{Classe A}
\[
1;\ 2;\ 2;\ 2;\ 2;\ 2;\ 3;\ 3;\ 3;\ 3;\ 3;\ 3;\ 3;\ 3;\ 6;\ 6;\ 6;\ 6;\ 6;\ 7;\ 7;\ 7
\]

\medskip
\textbf{Classe B}

\begin{center}
\begin{tabular}{>{\bfseries}l c}
Effectif total & 25\\
Moyenne & 4\\
Médiane & 5
\end{tabular}
\end{center}
}

\begin{enumerate}
  \item Donne les résultats de la classe A sous la même forme que ceux de la classe B.

  \item On appelle \textbf{grand lecteur} un élève qui a emprunté au moins \(5\) livres.

  \textbf{Quelle classe a le plus de grands lecteurs ? Justifie.}

  \item Peut-on connaître exactement le nombre de grands lecteurs de la classe B avec les seules informations données ? Explique.
\end{enumerate}

\vspace{0.5cm}

\begin{center}
  {\sffamily\LARGE\bfseries\color{TitreC}Séquence 1 -- Statistiques}\\[0.2em]
  {\sffamily\large\bfseries\color{SavoirC}Activité 4 -- Les grands lecteurs}
\end{center}

Deux classes du collège ont répondu à la question :

\begin{center}
\textit{« Combien de livres avez-vous empruntés durant les douze derniers mois ? »}
\end{center}

\blocdonnees{
\textbf{Classe A}
\[
1;\ 2;\ 2;\ 2;\ 2;\ 2;\ 3;\ 3;\ 3;\ 3;\ 3;\ 3;\ 3;\ 3;\ 6;\ 6;\ 6;\ 6;\ 6;\ 7;\ 7;\ 7
\]

\medskip
\textbf{Classe B}

\begin{center}
\begin{tabular}{>{\bfseries}l c}
Effectif total & 25\\
Moyenne & 4\\
Médiane & 5
\end{tabular}
\end{center}
}

\begin{enumerate}
  \item Donne les résultats de la classe A sous la même forme que ceux de la classe B.

  \item On appelle \textbf{grand lecteur} un élève qui a emprunté au moins \(5\) livres.

  \textbf{Quelle classe a le plus de grands lecteurs ? Justifie.}

  \item Peut-on connaître exactement le nombre de grands lecteurs de la classe B avec les seules informations données ? Explique.
\end{enumerate}

% NOTES PROFESSEUR -- NON IMPRIMÉES
% Adaptation de la tâche Éduscol « Les grands lecteurs » ; étendue retirée.
% Classe A : effectif 22 ; moyenne 86/22 ≈ 3,9 ; médiane 3 ; grands lecteurs 8.
% Classe B : médiane 5 et effectif 25 => au moins 13 valeurs >= 5.
% Donc B a nécessairement plus de grands lecteurs que A, sans connaître le nombre exact.

\end{document}
